\chapter{Counting to 20}\label{ch:g1:counting}

How many marbles? Counting answers the very first question of
mathematics. In this chapter we count small collections, write the
numbers from $0$ to $20$, and compare: who has more?

\section{Counting objects}

\begin{method}[How to count]\label{met:g1:counting:how}
\begin{enumerate}
  \item Touch each object once, saying the numbers in order: one,
        two, three, \dots;
  \item do not skip any object, do not touch one twice (move the
        counted ones aside, or cross them out);
  \item the \emph{last} number said is how many there are.
\end{enumerate}
The order of touching does not matter: counted left to right or
right to left, $7$ marbles are $7$ marbles.
\end{method}

\begin{omfigure}
\begin{tikzpicture}[scale=0.75]
  \foreach \i/\n in {0/1,1/2,2/3,3/4,4/5} {
    \fill[omDef] (\i*1.1,0) circle (0.33);
    \node[above, font=\small] at (\i*1.1,0.45) {\n};
  }
  \node[right, font=\small] at (5.2,0) {five marbles: there are $5$};
\end{tikzpicture}

{\small Counting: one number for each marble, and the last number is
the answer.}
\end{omfigure}

\begin{definition}[Zero]\label{def:g1:counting:zero}
When there is nothing to count --- an empty box --- the number is
\emph{zero}\index{zero}, written $0$.
\end{definition}

\section{The numbers up to 20}

\begin{example}[Reading and writing]\label{ex:g1:counting:writing}
The numbers from $0$ to $20$, in order:
\[
0,\ 1,\ 2,\ 3,\ 4,\ 5,\ 6,\ 7,\ 8,\ 9,\ 10,
\]
\[
11,\ 12,\ 13,\ 14,\ 15,\ 16,\ 17,\ 18,\ 19,\ 20 .
\]
After $9$, the numbers use \emph{two} digits: $13$ is a ten and $3$
more; $20$ is two tens. (Tens take the stage in
\cref{ch:g1:tens}.)
\end{example}

\begin{omfigure}
\begin{tikzpicture}[scale=0.62]
  \draw[-Stealth] (-0.4,0) -- (20.8,0);
  \foreach \i in {0,...,20} {
    \draw (\i,-0.12) -- (\i,0.12);
    \node[below, font=\footnotesize] at (\i,-0.15) {\i};
  }
  \fill[omThm] (7,0) circle (2.6pt);
  \fill[omMeth] (13,0) circle (2.6pt);
\end{tikzpicture}

{\small The number line from $0$ to $20$: every number has its place.
Walking right, the numbers grow. The dots mark $7$ and $13$.}
\end{omfigure}

\begin{example}[Just before, just after]\label{ex:g1:counting:neighbors}
On the line, each number has two neighbors: just before $14$ comes
$13$, just after comes $15$. Counting backward --- $10$, $9$, $8$,
\dots\ --- is walking left.
\end{example}

\section{Comparing}

\begin{method}[Who has more?]\label{met:g1:counting:compare}
To compare two collections, make pairs: one from each collection,
side by side.
\begin{enumerate}
  \item If some objects of one collection are left without a partner,
        that collection has \emph{more};
  \item if nothing is left over on either side, there are
        \emph{as many} on both.
\end{enumerate}
With numbers, no pairing is needed: the one further right on the
number line is the bigger --- pairing and the line always agree. We
write $12 > 9$ ($12$ is more) or $9 < 12$ ($9$ is less).
\end{method}

\begin{omfigure}
\begin{tikzpicture}[scale=0.62]
  \foreach \i in {0,...,6} \fill[omDef] (\i*1.1,1.3) circle (0.3);
  \foreach \i in {0,...,4} \fill[omThm] (\i*1.1,0) circle (0.3);
  \foreach \i in {0,...,4}
    \draw[omExo, thick] (\i*1.1,0.3) -- (\i*1.1,1.0);
  \draw[omExo, thick] (5.5,1.3) circle (0.52);
  \draw[omExo, thick] (6.6,1.3) circle (0.52);
  \node[right, font=\small] at (7.5,0.65)
    {$7 > 5$: two blue marbles have no partner};
\end{tikzpicture}

{\small Comparing by pairing: $7$ blue marbles against $5$ red ones
--- two blues (circled) are left over, so $7$ is more than $5$.}
\end{omfigure}

\begin{example}[First, second, third]\label{ex:g1:counting:ordinal}
In a race, the runners arrive in order: the \emph{first}, the
\emph{second}, the \emph{third}, the \emph{fourth}\,\dots\ These
words tell the \emph{position}, not the amount: the third runner is
one runner, at place number $3$.
\end{example}

\section{Exercises}

\begin{exercise}[$\star$]\label{exo:g1:counting:1}
Count the letters of your first name. Count the fingers of two
hands. Count the legs of two cats.
\end{exercise}

\begin{exercise}[$\star$]\label{exo:g1:counting:2}
Write in digits: seven; \ twelve; \ twenty; \ \omterm{def:g1:counting:zero}{zero}. Write in words:
$4$; \ $15$; \ $18$.
\end{exercise}

\begin{exercise}[$\star$]\label{exo:g1:counting:3}
Continue the counting: $8,\ 9,\ 10,\ ?,\ ?,\ ?$. Count backward:
$15,\ 14,\ 13,\ ?,\ ?,\ ?$.
\end{exercise}

\begin{exercise}[$\star$]\label{exo:g1:counting:4}
Which number comes just before $10$? Just after $16$? Between
$12$ and $14$?
\end{exercise}

\begin{exercise}[$\star$]\label{exo:g1:counting:5}
Copy and complete with $<$ or $>$:
\[
6 \;?\; 9, \qquad
14 \;?\; 8, \qquad
11 \;?\; 13, \qquad
20 \;?\; 19 .
\]
\end{exercise}

\begin{exercise}[$\star$]\label{exo:g1:counting:6}
Draw $6$ circles and $8$ crosses. Pair them up: which are more, and
by how many?
\end{exercise}

\begin{exercise}[$\star$]\label{exo:g1:counting:7}
Order from smallest to biggest: $12$; \ $7$; \ $17$; \ $2$; \ $10$.
\end{exercise}

\begin{exercise}[$\star$]\label{exo:g1:counting:8}
Five children stand in line: Ana, Bob, Cleo, Dan, Emma. Who is
second? Who is fourth? At which place is Dan?
\end{exercise}

\begin{exercise}[$\star\star$]\label{exo:g1:counting:9}
I am a number between $10$ and $20$. Counting by the line, I am
three steps after $13$. Who am I? And who is three steps
\emph{before} $13$?
\end{exercise}
